The registered title reflects the robotic welding proposal, and the implementation treats physics-informed fusion as conditioning on observed physical process context without imposing differential-equation or conservation-law constraints. After an audit documented development exposure to the proposed Intel Welding test recordings, evaluation centred on five adjacent industrial and robotic tasks, including a visual comparison of the frozen ResNet18-vd CNN backbone in RT-DETR-R18vd alone and with its hybrid Transformer encoder. A review of 94 full-text works examined fusion methods, single-input baselines and tests of missing or degraded inputs. In one weld-penetration study, images, sound and their combination achieved identical clean-test accuracy.

Independent validation of weld quality remains open. In PHM, already under way when the welding split was audited, three sensor groups predicted a photograph-derived gear-condition measure with a mean absolute error of 0.824 percentage points, compared with 0.804 for the training mean. Across 4,530 completed REASSEMBLE actions from 148 recordings, failure average precision reached .429 for sensors alone, .446 for equal averaging of video and sensor predictions, and .447 for logistic stacking. The difference between fusion methods remained uncertain. Two adaptive gates showed no consistent advantage under input loss or degradation. Sensor models outperformed stacking in MSM milling and ViFailback manipulation, and equal averaging in human–robot handover. Assembly provided the clearest evidence of a fusion benefit. Across tasks, the results support evaluating learned fusion against both a strong sensor model and simple averaging.
